\subsection{Collecting Data from Heterogeneous Sources}\label{subsec:big-data--collecting-data-from-heterogeneous-sources}

A modern data system rarely receives information from a single, clean, uniform source. Instead, it may combine data from sources such as WiFi logs, telecommunication systems, transport systems, IoT devices, social media, and web/mobile applications. The key word is: \textbf{heterogeneous}, meaning that the \hl{sources differ in structure, format, semantics, update frequency, and quality.}

\highspace
\begin{flushleft}
    \textcolor{Green3}{\faIcon{random} \textbf{Why heterogeneity matters}}
\end{flushleft}
Suppose a city wants to understand mobility. It might collect bus GPS data, metro ticket validations, WiFi connection logs, mobile network data, traffic sensors, and social media posts. All of these may provide information about movement, but they are \textbf{very different technically}. A GPS record might look like:
\begin{lstlisting}
vehicle_id = 51
lat = 45.4642
lon = 9.1900
timestamp = 17:31:08
\end{lstlisting}
A WiFi log might instead contain:
\begin{lstlisting}
device_hash
access_point
connection_time
disconnect_time
\end{lstlisting}
A social media post might be mostly text:
\begin{center}
    ``\emph{Traffic is terrible near Paris today}''
\end{center}
So \hl{one analytical problem may require data from several incompatible representations.}

\highspace
\begin{flushleft}
    \textcolor{Green3}{\faIcon{layer-group} \textbf{Different kinds of heterogeneity}}
\end{flushleft}
There are several dimensions of heterogeneity.
\begin{itemize}
    \item \definition{Structural Heterogeneity}: the data may have different structures (e.g., relational tables, JSON documents, text, images, sensor streams). This already creates a \hl{data-management challenge.} A relational schema such as:
\begin{lstlisting}
Customer(id, name, age)
\end{lstlisting}
        Is very different from a JSON event:
\begin{lstlisting}
{
  "device": "sensor17",
  "temperature": 24.8,
  "timestamp": "..."
}
\end{lstlisting}
        And both differ completely from an image.

    \item \definition{Semantic Heterogeneity}: even when two datasets have similar structure, the \textbf{meaning} may differ. Suppose one system stores \texttt{customer\_id} and another stores \texttt{user\_id}. \emph{Are they the same entity?} Maybe yes, maybe no. Or one dataset stores \texttt{location = "Milano"} while another uses \texttt{city\_code = 015146}. The problem is no longer syntax; it is: \hl{\emph{do these values represent the same real-world concept?}} This is a classic \textbf{data integration problem}.

    \item \definition{Temporal Heterogeneity}: different systems produce data at different frequencies. For example, a weather station may update every 10 minutes, while GPS data may update every second, and sales reports may be generated once per day. When combining sources, we must ask: \emph{Are these observations temporally comparable?} A traffic reading from 17:00 and a weather reading from 12:00 may not be meaningful together.

    \item \definition{Quality Heterogeneity}: different sources have different reliability. A calibrated sensor may be precise, while social media may contain false information. When combining sources, we must ask: \emph{How much can we trust each source?} This anticipates the later concept of \textbf{Veracity}.
\end{itemize}

\begin{flushleft}
    \textcolor{Green3}{\faIcon{cogs} \textbf{Collection is therefore more than copying bytes}}
\end{flushleft}
A simplistic (and naive) picture would be:
\begin{equation*}
    \text{source}
    \rightarrow
    \text{database}
\end{equation*}
But in a realistic data system, ingestion often looks more like:
\begin{enumerate}
    \item Sources
    \item Acquisition
    \item Parsing
    \item Cleaning
    \item Transformation
    \item Integration
    \item Storage
\end{enumerate}
Not every pipeline performs all these steps immediately, but these are the kinds of issues involved.

\highspace
\begin{examplebox}[: Combining transport and telecom data]
    Suppose we want to estimate how many people move between two areas of a city. We could combine \emph{public transport validations} with \emph{mobile network presence} and perhaps \emph{WiFi connections}.

    \highspace
    Each source gives a different approximation of human movement. But before combining them, we must solve questions like:
    \begin{itemize}
        \item \emph{What geographic granularity does each source use?} (e.g., GPS coordinates, cell tower coverage, bus stop locations)
        \item \emph{How are timestamps represented?} (e.g., local time, UTC, different time zones)
        \item \emph{How do we avoid counting the same person twice?} (e.g., if a person uses both public transport and their mobile phone)
        \item \emph{Do we have personally identifiable information?} (e.g., phone numbers, device IDs, names)
        \item \emph{Are identifiers stable?} (e.g., do device IDs change over time, do people change phone numbers)
        \item \emph{Are some observations missing?} (e.g., if a person doesn't use public transport or doesn't carry a phone)
    \end{itemize}
    So heterogeneous-data collection is \textbf{not just a storage problem}. It is also an \textbf{integration problem}.
\end{examplebox}

\highspace
\begin{deepeningbox}[: Cross-source integration]
    This concept will become important later. Suppose we have:
    \begin{equation*}
        S_1,S_2,\ldots,S_n
    \end{equation*}
    Different sources. The \textbf{goal} is not necessarily just to store all of them independently. Often we want:
    \begin{equation*}
        S_1 + S_2 + \cdots + S_n
        \rightarrow
        \text{integrated view}
    \end{equation*}
    Because the \hl{integrated dataset may reveal something that no source can reveal alone.}

    \highspace
    For example:
    \begin{equation*}
        \text{purchase history}
        +
        \text{web browsing}
        +
        \text{customer profile}
    \end{equation*}
    May support much better recommendation models than any one source independently. This is why \hl{combining sources creates potential value.}
\end{deepeningbox}

\highspace
\begin{takeawaysbox}
    The central point is that \hl{Big Data systems must often combine many heterogeneous data sources.} Those sources differ in structure, semantics, timing, and quality. Therefore ingestion is not merely \emph{moving data into storage}, but often also involves parsing, cleaning, transforming, aligning, and integrating the data. \hl{The payoff is that combining multiple sources can produce richer information than any one source alone.}
\end{takeawaysbox}